% Generated by tools/textbook/export_listings.js from tools/decks/mastery.js and tools/concepts/registry.json: do not edit.
\begin{chapstatus}{Stato di apprendimento di questo capitolo}
\masterychip{now}{DA COMPRENDERE OGGI}\enspace gli ingressi con pull-up; tensione, corrente, legge di Ohm; il transistor come interruttore; dai volt ai bit; le porte CMOS.\par\smallskip
\masterychip{use}{DA UTILIZZARE OGGI}\enspace il programma completo, gli strumenti di PRG32 (si usa, non si studia, in questo corso).\par\smallskip
\masterychip{later}{DA COMPRENDERE PIÙ AVANTI}\enspace il clock (Capitolo~\ref{ch:sequential}).
\end{chapstatus}
\begin{chapstatus}{Percorso del capitolo}
Alla fine del capitolo saprete:
\begin{enumerate}[nosep,leftmargin=2em]
\item spiegare come un transistor MOS funziona da interruttore e come CMOS costruisce una porta
\item leggere livelli logici e margini di rumore dalle tensioni
\item dimensionare la resistenza di un LED e spiegare il pull-up di un pulsante
\end{enumerate}
Si appoggia su: solo l'installazione del Capitolo~\ref{ch:setup}. Prossimo passo: Capitolo~\ref{ch:combinational}.
\par\smallskip\textbf{Guardando indietro, guardando avanti.} Nella Lezione 02 il contatore sopravviveva da un frame all'altro in memoria, non in un registro. Oggi una cella di memoria costa circa 6 transistor e un flip-flop circa 20: ricordare un bit è un circuito. La Lezione 13 mostra come latch e flip-flop conservano un bit. \emph{Richiamo:} Dove abitava il contatore della Lezione 02 tra due frame, e perché non in t1?
\end{chapstatus}
